%% notes-tex/010-additive-baselines/sections/6-results.tex

\section{Held-out results\secstatus{todo}}
\label{sec:results}

%% SOURCE: experiments/010-kuzmin-tmi/scripts/additive_baseline_gene_interaction.py
%% -> results/additive_baseline_gene_interaction.csv. Transformer rows are read from
%% the three checkpoints' re-evaluation runs under $DATA_ROOT/wandb-experiments.
\begin{table}[htbp]
\centering
\begin{tabular}{llrrr}
\hline
Model & Interaction capacity & Pearson & Spearman & MSE \\
\hline
B0 train mean & none & 0.000 & 0.000 & 0.004124 \\
B4 query pair only & pairwise, 420 pairs & 0.390 & 0.362 & 0.003493 \\
B3 hierarchical mean & pairwise, empirical & 0.390 & 0.362 & 0.003493 \\
B1 additive per-gene ridge & none & 0.400 & 0.406 & 0.003460 \\
B2 additive plus pair ridge & pairwise, 420 pairs & 0.406 & 0.412 & 0.003442 \\
B5 embedding MLP, seed 0 & full & 0.425 & 0.417 & 0.003382 \\
B5 embedding MLP, seed 1 & full & 0.427 & 0.421 & 0.003395 \\
B5 embedding MLP, seed 2 & full & 0.425 & 0.419 & 0.003381 \\
CGT M01 \texttt{lzs9pcj3} & full & 0.443 & 0.426 & 0.003400 \\
CGT M02 \texttt{yv4r30bi} & full & 0.438 & 0.424 & 0.003439 \\
CGT M03 \texttt{c7671wgj} & full & 0.455 & 0.426 & 0.003315 \\
\hline
\end{tabular}
\caption{Test split, 37{,}673 records, raw $\tau$ units. The interaction-capacity
column repeats Table~\ref{tab:capacity}: only the last four rows can express a
three-way interaction, which is what the label is. Ridge penalties are selected
on validation; B5 early-stops on validation; the transformer checkpoints are
selected on validation. Test was not consulted by any of them.}
\label{tab:heldout}
\end{table}

\tcfigfit[0.72]{figures/additive_baseline_test_pearson.pdf}{Test Pearson for
every model. Error bars on the nonlinear baseline are the standard deviation
across three seeds; the transformer bars are three separate training runs shown
individually.}{fig:ladder}

\subsection{Reading the ladder\secstatus{todo}}

\notesec{The additive null is close, but the transformer does beat it}
B1 reaches 0.400 against the transformer's 0.438 to 0.455, so a model that
cannot represent any interaction recovers 88 to 91 percent of the correlation.
In explained variance the gap is $r^2 = 0.160$ against $0.207$ for the best
checkpoint, an increment of 0.047 of label variance. In error, the best
checkpoint reduces mean squared error 19.6 percent below B0 and B1 already
delivers 16.1 percent of that. This is not the DeWitt result, where the network
collapsed onto the additive model at $r > 0.999$. The margin here is real. It is
also small.

\notesec{Seed spread is a large fraction of that margin}
The three transformer runs span 0.438 to 0.455, a standard deviation of 0.0088.
The three B5 seeds span 0.425 to 0.427, a standard deviation of 0.0011. So the
transformer's run-to-run spread is eight times the nonlinear baseline's and is
roughly a fifth of its own margin over B1. A single checkpoint's number is not
the capability.

\notesec{Capacity, not the graphs, closes most of the remaining gap}
B5 sees exactly B1's features, has no graph and no attention, and reaches 0.426.
That closes about half the distance from B1 to the transformer. On mean squared
error B5's 0.003382 is better than two of the three checkpoints. Since
Section~\ref{sec:graphs} establishes that the nine graphs do not enter the
forward pass, none of the B5-to-transformer difference can be attributed to
them; what is left is the richer set aggregation of Eq.~\ref{eq:cgt} and roughly
eight times the parameters.

\notesec{B4 is the result that changes what the metric means}
B4 knows only which of 420 recurring gene pairs a triple contains. It never
looks at the third gene. It reaches 0.390 test Pearson, 97 percent of B1's and
86 percent of the best checkpoint's. Those 420 pairs are the Kuzmin query
doubles, the most frequent appearing 3{,}308 times, so B4 is close to a
per-screen offset.

On a split that is random over records, then, most of what every model in
Table~\ref{tab:heldout} is scoring is the identity of the screen a triple came
from. That is not trigenic biology and it would not survive a query-pair-disjoint
split. It also explains why the additive null does so well on a label that was
constructed to have the additive part removed: what B1's per-gene coefficients
absorb is largely query-pair structure reaching them through the genes that make
up the query pairs.

\notesec{Spearman compresses the ladder further}
The rank correlations run 0.406 for B1, 0.412 for B2 and 0.417 to 0.421 for B5,
against 0.424 to 0.426 for the transformer. The gap from the additive null to
the transformer is 0.020 in Spearman against 0.045 in Pearson. Ranking is what
a design-space selection consumes, so the smaller of the two is the more
decision-relevant number.

\subsection{Verifying the architecture reading\secstatus{todo}}
\label{sec:verify}

Section~\ref{sec:sees} claims the encoder is strain-independent and the nine
graphs do not enter the forward pass. Both are checkable rather than merely
readable, and checking them also produces the per-record predictions that
Section~\ref{sec:limits} lists as missing.

The stock evaluation path cannot run: the 010 LMDB stores
\texttt{perturbation\_type: "deletion"}, which the current schema rejects. The
route around it follows from the architecture. If the encoder really does not
see the strain, then scoring a record needs only the perturbed gene indices, and
the index space is the sorted base-graph node list, which is rebuildable from
the genome and the build's \texttt{gene\_set.json} without touching the LMDB.
\texttt{score\_010\_checkpoints\_directly.py} does that, loads each checkpoint
into the model class, and recomputes the metrics.

The check is that the recomputed numbers must reproduce what the original
evaluation runs recorded. A wrong gene ordering or a mis-loaded weight would
give a plausible but wrong number, so agreement to several digits is the
evidence that the reconstruction is faithful.

%% SOURCE: experiments/010-kuzmin-tmi/scripts/score_010_checkpoints_directly.py
%% -> results/cgt_direct_scoring.csv, against the recorded eval-run summaries
\begin{table}[htbp]
\centering
\begin{tabular}{lrrrr}
\hline
& \multicolumn{2}{c}{Val Pearson} & \multicolumn{2}{c}{Val MSE} \\
Checkpoint & recomputed & recorded & recomputed & recorded \\
\hline
M01 \texttt{lzs9pcj3} & 0.451970 & 0.451963 & 0.003244 & 0.003244 \\
M02 \texttt{yv4r30bi} & 0.447319 & 0.447155 & 0.003258 & 0.003259 \\
M03 \texttt{c7671wgj} & 0.461917 & 0.461881 & 0.003163 & 0.003163 \\
\hline
\end{tabular}
\caption{Scoring the checkpoints through a hand-built forward pass that computes
the encoder once and feeds it only perturbed-gene indices reproduces the
recorded metrics. The residual difference is at the fourth decimal and is
consistent with the original runs using mixed bf16 precision.}
\label{tab:repro}
\end{table}

Two details had to be right for this to work, and both are worth recording
because getting either wrong produced a plausible but wrong answer. The index
space is the sorted genome gene set, 6{,}607 genes, not the build's own 6{,}579
gene set; using the latter gave a validation Pearson of $-0.001$ to $0.002$. And
the readout pools by \emph{mean}, whereas the current class defaults to
\texttt{sum}; using the default gave 0.420 to 0.438, close enough to look
credible and wrong.

That reproduction is the direct verification of Section~\ref{sec:encoder}. The
encoder was evaluated exactly once, on a strain-independent input, and the
recorded per-split metrics came back. If the encoder had needed the strain, this
could not have happened.

The graph ablation gives the other claim. Setting the graph penalty weight from
$1.0$ to $0.0$ changes the maximum absolute test prediction by
$9.0 \times 10^{-7}$, $1.2 \times 10^{-6}$ and $1.3 \times 10^{-6}$ for the three
checkpoints, against a prediction spread of about $0.036$. The nine graphs do not
enter the forward pass.

\subsection{Record-by-record agreement\secstatus{todo}}
\label{sec:paired}

With per-record predictions available, the comparison can be made the way the
DeWitt preprint makes it, between the predictions themselves rather than between
their summary statistics.

%% SOURCE: experiments/010-kuzmin-tmi/scripts/paired_prediction_agreement.py
%% -> results/paired_prediction_agreement.csv
\begin{table}[htbp]
\centering
\begin{tabular}{lrrr}
\hline
& \multicolumn{3}{c}{Pearson with} \\
Model & CGT M01 & CGT M02 & CGT M03 \\
\hline
B1 additive ridge & 0.752 & 0.726 & 0.759 \\
B2 additive plus pair & 0.759 & 0.734 & 0.765 \\
B4 query pair only & 0.729 & 0.706 & 0.735 \\
B5 embedding MLP & 0.788 & 0.784 & 0.810 \\
CGT M01 & & 0.788 & 0.821 \\
CGT M02 & 0.788 & & 0.804 \\
\hline
\end{tabular}
\caption{Correlation between models' test-set predictions, 37{,}673 records.
The transformer does not reproduce the additive model, which is where 010
departs from the MULTI-evolve case. But two transformer training runs agree with
each other at 0.788 to 0.821, barely more than a transformer agrees with the
additive model at 0.726 to 0.759, and no more than it agrees with the nonlinear
baseline at 0.784 to 0.810.}
\label{tab:paired}
\end{table}

\tcfigfit[0.98]{figures/paired_prediction_agreement.pdf}{Test-set predictions of
one model against another. The transformer is no closer to a second training run
of itself than it is to a nonlinear baseline that has no graph and an eighth of
the parameters.}{fig:paired}

\notesec{The headline distinction from the preprint}
DeWitt's networks correlated with the additive model at $r > 0.999$, so the
network \emph{was} the additive model. Here the correlation is $0.73$ to $0.76$,
so the 010 transformer is genuinely computing something the additive model does
not. That conclusion should be stated plainly, because it is the opposite of the
preprint's finding and it is what the null was fit to test.

\notesec{What that extra content is worth is a separate question}
Seed-to-seed correlation between two transformer runs is $0.79$ to $0.82$. The
part of the transformer's output that is not additive is therefore largely not
reproducible across training runs: a checkpoint disagrees with the additive
model about as much as it disagrees with another checkpoint of the same
architecture on the same split.

\notesec{The models are wrong in the same places}
Correlating the residuals $y - \hat{y}$ of B1 and each checkpoint gives $0.907$,
$0.897$ and $0.918$. Almost all of the error is shared, so the transformer is
not fixing a class of records the additive model gets wrong; it is making a
modest, broadly distributed improvement.

\notesec{The advantage is real and bounded}
A paired bootstrap over test records, 2{,}000 resamples, puts the Pearson gain
over B1 at $+0.043$ with a 95 percent interval $[+0.023, +0.063]$ for M01,
$+0.038$ with $[+0.019, +0.058]$ for M02, and $+0.055$ with $[+0.034, +0.074]$
for M03. All three intervals exclude zero. B5 gains $+0.025$ with
$[+0.017, +0.033]$. So the transformer beats a no-interaction null by a margin
that is statistically established and is between two and six hundredths of a
correlation.

\notesec{Rankings diverge where selection happens}
Among the 100 records each model calls most negative, B1 and M03 share 31,
B5 and M03 share 45, and B1 and B5 share 46. At $K = 10{,}000$ all pairs exceed
0.91. The models agree about the bulk and disagree about the tail, which is the
part a design-space selection reads.

\notesec{Prediction spread}
Test prediction standard deviations are 0.025 for B1, 0.029 for B5 and 0.036 to
0.038 for the checkpoints, against a test label standard deviation of 0.064.
Every model is heavily shrunk toward the mean, as a squared-error fit on a noisy
target should be.
